\section{The Moving-Knife Algorithm}
\label{sec:algorithm}

\subsection{Notation}

Let $G_v=(V_v,E_v)$ be the current bisection-tree node with $m=|V_v|$ tracts.
Each tract $i$ has population $p_i$ and centroid
$\ell_i=(\lambda_i,\phi_i)$. The implementation projects centroids to Albers
coordinates $\mathbf{x}_i=(x_i,y_i)$ before scoring directions.

\subsection{The Implemented Split}

For each orientation $\theta_k=k\pi/n_{\mathrm{orient}}$, \mka{} computes the
projection
\[
  q_i(\theta_k) = x_i\cos\theta_k + y_i\sin\theta_k.
\]
Sorting by $q_i$ turns the node into a one-dimensional queue. The split point
is the first prefix whose cumulative population reaches $P_v/2$.

\begin{center}
\begin{tabular}{r|rrrrrrrr|l}
rank & 1 & 2 & 3 & 4 & 5 & 6 & 7 & 8 & candidate \\
\hline
pop  & 4 & 3 & 2 & 6 & 1 & 5 & 2 & 3 & $P_v=26$ \\
cum  & 4 & 7 & 9 & 15 & 16 & 21 & 23 & 26 & stop at 15 \\
side & L & L & L & L & R & R & R & R & $L/R$
\end{tabular}
\end{center}

The score for this candidate is
\[
  \min(\reock{}_{\mathrm{centroid}}(L),\reock{}_{\mathrm{centroid}}(R)).
\]
The algorithm keeps the assignment with the largest score across all tested
orientations.

\subsection{Seed Derivation}
\label{subsec:seed}

The code defines a deterministic node seed:
\[
  s_v^{\mathrm{mka}} = \mathrm{first8}\bigl(
    \mathrm{SHA\text{-}256}(
      \texttt{"MKA\_INIT\_"}
      \;\|\; \mathrm{le4}(|\mathtt{node\_path}|)
      \;\|\; \mathtt{node\_path}
      \;\|\; \mathrm{le8}(\mathtt{base\_seed})
    )
  \bigr).
\]
The prefix is tested to be distinct from CVD seed prefixes. In the current
split routine this value is computed as a reserved tie seed, but ties are
resolved by the first strictly best orientation encountered rather than by a
random draw. The seed formula should still be retained in manifests once
per-node angle ledgers are archived.

\subsection{Pseudocode}

\begin{algorithm}[t]
\caption{\mka{}-\textsc{Bisect}($G_v$, $\mathbf{x}$, $\mathbf{p}$,
  $n_{\mathrm{orient}}$, $\tau$)}
\label{alg:mka}
\begin{algorithmic}[1]
\State sort $V_v$ by global tract id for deterministic local indexing
\State build local adjacency for the later boundary rebalance
\State $P \gets \sum_{i\in V_v}p_i$;\quad $P^* \gets \lfloor P/2\rfloor$
\State project all available centroids to Albers coordinates
\State $\mathtt{best\_score}\gets-\infty$;\quad
       $\mathtt{best\_assignment}\gets$ all right
\For{$k=0$ \textbf{to} $n_{\mathrm{orient}}-1$}
  \State $\theta\gets k\pi/n_{\mathrm{orient}}$
  \State compute $q_i=x_i\cos\theta+y_i\sin\theta$ for every local tract $i$
  \State sort tracts by $q_i$
  \State scan the sorted list until cumulative population reaches $P^*$
  \State assign the prefix left and the suffix right
  \State compute centroid-MEC Reock for left and right
  \State $\mathtt{score}\gets\min(r_L,r_R)$
  \If{$\mathtt{score}>\mathtt{best\_score}$}
    \State keep this assignment
  \EndIf
\EndFor
\State rebalance for at most 200 boundary moves toward $P/2$
\State repair degenerate empty-side cases if necessary
\State \Return left and right global tract sets
\end{algorithmic}
\end{algorithm}

\subsection{Boundary Rebalance}

The sweep can overshoot the half-population target by one tract. After the best
orientation is chosen, the code repeatedly moves a boundary tract from the
heavier side to the lighter side when that move best reduces the population
excess. This loop is capped at 200 iterations.

The loop is deliberately local. It considers only tracts adjacent to the other
side, but it does not run a full contiguity proof after every tentative move.
As with the other structure splitters, the returned plan should be validated by
the downstream plan checks before any certification claim is made.

\subsection{Data and CLI}
\label{subsec:data}

\mka{} requires tract centroids. If centroids are absent, the CLI returns a
configuration error:
\begin{verbatim}
[CONFIG] --structure moving-knife requires centroid data.
Run: bisect fetch --type centroids
\end{verbatim}

The standard invocation is:
\begin{verbatim}
bisect state --state NC --year 2020 \
  --structure moving-knife \
  --mka-orientations 180 \
  --mka-metric reock
\end{verbatim}

\texttt{--mka-metric polsby-popper} is accepted by the enum, but currently
falls back to the Reock scoring path in \texttt{split\_subgraph\_mka}.

\subsection{Runtime}

Each orientation requires $O(m)$ projection work, $O(m\log m)$ sorting, a
linear population scan, and two minimum-enclosing-circle scores. The resulting
node-level runtime is
\[
  O(n_{\mathrm{orient}}\,m\log m + 200m).
\]
For the default $n_{\mathrm{orient}}=180$, this is heavier than one simple
graph split but still small enough to serve as an auditable geometric baseline.
